% Version: v47 (2026-10-16), every result from another reference is stated with its citation before use; locators and statements corrected against the sources. No claim changes.
\documentclass[]{interact}
\usepackage[utf8]{inputenc}

\usepackage{mathtools}
\usepackage{amsmath,amssymb,amsthm}
\usepackage{amsfonts}
\usepackage{enumerate}
\usepackage{enumitem}

\usepackage[hidelinks,pagebackref=false]{hyperref}
\usepackage[all]{hypcap}
\hypersetup{linktoc=all}

\usepackage{tikz}
\usepackage[yyyymmdd,hhmmss]{datetime}
\usepackage{float}

\usepackage[numbers,sort&compress]{natbib}
\bibpunct[, ]{[}{]}{,}{n}{,}{,}
\renewcommand\bibfont{\fontsize{10}{12}\selectfont}
\makeatletter
\def\NAT@def@citea{\def\@citea{\NAT@separator}}
\makeatother

\theoremstyle{plain}
\newtheorem{theorem}{Theorem}[section]
\newtheorem{lemma}[theorem]{Lemma}
\newtheorem{corollary}[theorem]{Corollary}
\newtheorem{proposition}[theorem]{Proposition}

\theoremstyle{definition}
\newtheorem{definition}[theorem]{Definition}
\newtheorem{example}[theorem]{Example}

\theoremstyle{remark}
\newtheorem{remark}[theorem]{Remark}
\newtheorem{notation}{Notation}

\allowdisplaybreaks

\newcommand{\Der}{\mathrm{Der}}
\newcommand{\Inn}{\mathrm{Inn}}
\newcommand{\Hom}{\operatorname{Hom}}
\newcommand{\JDer}{\mathrm{JDer}}
\newcommand{\BSDer}{\mathrm{BSDer}}
\newcommand{\ADer}{\mathrm{Der}_\alpha}
\newcommand{\LDer}{\mathrm{LDer}}


\newcommand{\paperversion}{46}


\begin{document}
\title{Braided-Symmetric Derivations of Algebras}

\author{
\name{Abderrahim Adrabi\textsuperscript{a,b}\thanks{CONTACT A. Adrabi. Email: abderrahim.adrabi@um5r.ac.ma}}
	\affil{\textsuperscript{a}Faculty of Sciences, Mohammed V University in Rabat, Morocco;\\ \textsuperscript{b}Faculty of Educational Sciences, Mohammed V University in Rabat, Morocco}
}

\history{Preprint version \paperversion, compiled \today}

\maketitle

\begin{abstract}
We study braided-symmetric derivations of an algebra over a field.
These are maps that satisfy the Leibniz rule for the product symmetrized by an involution of $A \otimes A$.
The choice of involution decides which derivations appear.
The identity involution gives the ordinary derivations when $2$ is invertible.
The swap gives the Jordan derivations, and its antisymmetric version gives the Lie derivations.
So we ask when a Jordan or Lie derivation is a derivation.
This question becomes a comparison of two involutions on one algebra.
On full matrix algebras and on upper triangular $2 \times 2$ matrices, a trace-type map gives an example.
It is a Lie derivation but not a derivation.
We then twist the swap by an involutive automorphism.
Let the algebra be commutative and reduced, with $2$ invertible.
Then the twisted maps that commute with the automorphism are exactly the derivations that commute with it.
An example with nilpotent elements shows that the reduced hypothesis is needed.
\end{abstract}

\begin{keywords}
Braided-symmetric derivations; twisted Jordan derivations; triangular algebras; Jordan derivations; Lie derivations
\end{keywords}

\begin{amscode}
MSC 2020: 16E40; 18M05; 16W25
\end{amscode}

\section{Introduction}\label{sec:intro}

Throughout, $K$ is a field and an algebra means a unital associative $K$-algebra.
We take all maps between algebras to be $K$-linear.
We say so when a result cited below holds more generally for additive maps on rings.

A \emph{derivation} of an algebra $A$ is a linear map $D \colon A \to A$ with $D(xy) = D(x)y + xD(y)$ for all $x, y \in A$.
The notion goes back to \citet[Part I]{Jacobson1937}, who studied the derivations of an arbitrary nonassociative algebra.
\citet[\S\S2 and 4]{Hochschild1945} defined the cocycles of an associative algebra, and its $1$-cocycles are the derivations.
A \emph{Jordan derivation} of $A$ is a linear map with $D(x \circ y) = D(x) \circ y + x \circ D(y)$ for all $x, y \in A$, where $x \circ y = xy + yx$.
\citet[p.~1104]{Herstein1957} introduced this notion.
His result needs three notions for a ring $R$.
We call $R$ \emph{prime} if $aRb = 0$ implies $a = 0$ or $b = 0$ \citep[p.~1003]{Bresar1988}.
We call $R$ \emph{semiprime} if $aRa = 0$ implies $a = 0$ \citep[p.~1003]{Bresar1988}.
We call $R$ \emph{$2$-torsion free} if $2a = 0$ implies $a = 0$ \citep[p.~1003]{Bresar1988}.
So every prime ring is semiprime.
Herstein proved that every Jordan derivation of a prime ring of characteristic different from $2$ is a derivation \citep[Theorem 3.1]{Herstein1957}.
Every derivation is a Jordan derivation, and the converse is false in general \citep[p.~1003]{Bresar1988}.
\citet[Corollary 5]{Cusack1975} extended this result of Herstein to $2$-torsion free semiprime rings, and \citet[Theorem 1]{Bresar1988} gave an independent proof.

A \emph{Lie derivation} of $A$ is a linear map with $D([x,y]) = [D(x), y] + [x, D(y)]$, where $[x,y] = xy - yx$.
\citet[p.~183]{Martindale1964} studied the Lie derivations of a primitive ring of characteristic not $2$ that contains a nontrivial idempotent, and showed that each has the form $D + T$, with $D$ a derivation and $T$ a map into the centre that vanishes on commutators.
So for Lie derivations no result of Herstein type can hold.
Here is the reason.
Let $t \colon A \to Z(A)$ be a linear map into the centre with $t([x,y]) = 0$ for all $x, y$.
Every term of the Lie derivation identity for $t$ is then zero, so $t$ is a Lie derivation.
But $t$ is not a derivation as soon as $t(1) \neq 0$.

So a description of Lie derivations must allow such a term.
A Lie derivation has the \emph{standard form} if it equals $d + t$, with $d$ a derivation and $t$ a central-valued map that vanishes on commutators \citep[\S1]{Cheung2003}.
A ring is \emph{left primitive} if it has a faithful simple left module, and right primitive rings are defined in the same way \citep[Definition 11.2]{Lam2001}.
\citet[Theorem 2]{Martindale1964} proved that every Lie derivation has this form for primitive rings of characteristic not $2$ with a nontrivial idempotent.
Here an \emph{idempotent} is an element $e$ with $e^2 = e$, and it is \emph{nontrivial} if $e \neq 0$ and $e \neq 1$.
Here $d$ may take values in a larger ring.

The next result uses the \emph{standard polynomial} of degree $4$ \citep[\S1]{Amitsur1950},
\[
S_4(x_1, x_2, x_3, x_4) = \sum_{\varpi} \operatorname{sgn}(\varpi)\, x_{\varpi(1)} x_{\varpi(2)} x_{\varpi(3)} x_{\varpi(4)},
\]
where $\varpi$ runs over the permutations of $\{1,2,3,4\}$.
We say that a ring \emph{satisfies $S_4$} if $S_4$ is zero for every choice of four elements of the ring.
A prime ring satisfies $S_4$ exactly when it is commutative or is isomorphic to a subring of $M_2(F)$ for a field $F$ \citep[Lemma 1]{Bresar1993}.
\citet[Theorem 5]{Bresar1993} proved the standard form for prime rings of characteristic different from $2$ that do not satisfy $S_4$.

Triangular algebras are defined in Section~\ref{sec:preliminaries}.
For them, \citet[Theorem 2.1]{ZhangYu2006} proved that every Jordan derivation of a triangular algebra is a derivation.
For Lie derivations, \citet[Theorem 11]{Cheung2003} found conditions under which every Lie derivation has the standard form.
\citet[\S3]{LiWei2015} extended this to more algebras.
\citet[Theorem 3.1 and Corollary 3.1]{Wang2019} treated every triangular algebra, with no condition on the module or the centre. For these algebras a Lie derivation is a sum $\delta + \tau$, where $\delta$ is a derivation into a larger triangular algebra and $\tau$ is a linear map into its centre that vanishes on commutators.

In the result of Wang, every Lie derivation of a triangular algebra $\mathcal{U}$ is a sum $d + t$ \citep[Theorem 3.1]{Wang2019}.
Here $d$ is a derivation of $\mathcal{U}$ into a larger triangular algebra $\mathcal{U}_0$ with the same unit \citep[Theorem 3.1]{Wang2019}.
The map $t$ is linear, takes values in the centre of $\mathcal{U}_0$ and vanishes on commutators \citep[Theorem 3.1]{Wang2019}.

\citet[\S1]{Benkovic2015} treats derivations, Jordan derivations and Lie derivations as $f$-derivations, for $f = x_1x_2$, $x_1 \circ x_2$ and $[x_1, x_2]$, so that one identity for a multilinear polynomial $f$ covers all three.
We obtain them from an identity of a different kind.
Our identity keeps the polynomial $x_1x_2$, and an involution of $A \otimes A$ acts before we multiply.
Here the involution may use extra structure of $A$, such as an automorphism.
The identity is then no longer an $f$-derivation identity; Section~\ref{sec:preliminaries} explains why.

Let $A$ be an algebra over $K$ with product $\mu$.
Let $\sigma \colon A \otimes A \to A \otimes A$ be a $K$-linear map with $\sigma^2 = \mathrm{id}$, for instance the swap $\tau(x \otimes y) = y \otimes x$.
It gives a symmetrized product
\[
x \cdot_\sigma y = \mu(x \otimes y) + \mu(\sigma(x \otimes y)).
\]

For $\sigma = \tau$ this is the Jordan product $x \circ y$.
Put $\mu_\sigma = \mu + \mu \circ \sigma$.
Since $\sigma^2 = \mathrm{id}$, this gives
\[
\mu_\sigma \circ \sigma = \mu_\sigma.
\]
Example~\ref{ex:not-associative} shows that a twisted version of this product can fail to be associative and to have a unit.

We call a linear map $D \colon A \to A$ a \emph{braided-symmetric derivation} for $\sigma$ if
\begin{equation}\label{eq:sigma-leibniz}
D(x \cdot_\sigma y) = D(x) \cdot_\sigma y + x \cdot_\sigma D(y).
\end{equation}
So \eqref{eq:sigma-leibniz} is the Leibniz rule for the symmetrized product.
Other authors have studied \emph{braided derivations}, which use a braiding operator directly \citep[\S2]{Rozanski1996} \citep[\S4.1, Definition 14]{Huru2006} \citep[\S3, equation (3.9)]{Weber2020}.
Remark~\ref{rem:distinction-from-braided} compares the two identities.
Our identity uses only $\mu$, $\sigma$ and $D$.

We now describe the main results.
Our first step is a dictionary.
Remark~\ref{rem:sigma-derivation-classification} shows that the identity gives exactly the derivations, and the swap gives exactly the Jordan derivations.

Our second step twists the swap by an involutive automorphism or anti-automorphism $\alpha$ of $A$.
Proposition~\ref{prop:defect} describes the twisted class through the defect $E_D(x,y) = D(xy) - D(x)y - xD(y)$.
Corollary~\ref{cor:twisted-semiprime} shows that, on $2$-torsion-free semiprime rings, the Jordan maps in the twisted class are the derivations commuting with $\alpha$.

Theorem~\ref{thm:twisted-commutative} treats commutative reduced algebras and needs no Jordan hypothesis.
Example~\ref{ex:twisted-nilpotent} shows that the reduced hypothesis is needed.
Proposition~\ref{prop:twisted} shows that the twisted class and the $\alpha$-derivations on one algebra do not contain each other.
Proposition~\ref{prop:base-change} shows that these spaces keep their dimension when the base field grows.

Our third step uses the difference $\mu - \mu \circ \sigma$ in place of the sum.
Proposition~\ref{prop:lie-classification} shows that it gives the Lie derivations, together with the easy half of the theorem of Martindale and Bre\v{s}ar.
Proposition~\ref{prop:twisted-lie} describes the twisted version through the same defect.
Example~\ref{ex:matrix} finds both spaces for full matrix algebras.
Corollaries~\ref{cor:triangular-jordan} and~\ref{cor:triangular-lie} apply the results to triangular algebras, where the theorems of Zhang and Yu and of Cheung give complete answers.
Example~\ref{ex:t2} checks them on the smallest case.

We follow this order.
Section~\ref{sec:preliminaries} recalls the classical notions and contains no new results.
Section~\ref{sec:sigma-derivations} defines braided-symmetric derivations and treats the identity and the swap.
Section~\ref{sec:twisted} treats the twisted swap through the Leibniz defect.
Section~\ref{sec:twisted-commutative} treats reduced commutative algebras and compares the twisted class with $\alpha$-derivations.
Section~\ref{sec:lie} treats the antisymmetrized product.
Triangular algebras appear as corollaries in Sections~\ref{sec:sigma-derivations} and~\ref{sec:lie}.
We place each example right after the result it illustrates.

\section{Preliminaries}\label{sec:preliminaries}

This section recalls the classical notions we need, and none of it is new.
We treat five kinds of derivation and then triangular algebras.


From now on, $A$ is a $K$-algebra with unit $1$.
We write $[x,y] = xy - yx$ for the commutator, $x \circ y = xy + yx$ for the Jordan product, and $Z(A)$ for the centre.
The centre is the set of elements $z$ with $zx = xz$ for all $x \in A$.
We call an element $x$ \emph{nilpotent} if $x^n = 0$ for some $n \geq 1$.
A commutative algebra is \emph{reduced} if its only nilpotent element is zero \citep[1.4.12]{ChambertLoir2021}.

A derivation of an algebra is as defined by \citet[Part I]{Jacobson1937}: a linear map $D$ with $D(xy) = D(x)y + xD(y)$.
Every derivation satisfies $D(1) = 0$, since $D(1) = D(1 \cdot 1) = 2D(1)$.
For $g \in A$, the map $\operatorname{ad}_g(x) = [g,x]$ is an \emph{inner derivation}.
We write $\Der(A)$ and $\Inn(A)$ for these two spaces.

We add a word on the characteristic for Jordan derivations.
A \emph{Jordan derivation} in the sense of \citet[p.~1104]{Herstein1957} is an additive map with $D(x^2) = D(x)x + xD(x)$.
We write $\JDer(A)$ for the space of linear maps with $D(x \circ y) = D(x) \circ y + x \circ D(y)$ for all $x, y$.
Replacing $x$ by $x + y$ gives the condition $D(x \circ y) = D(x) \circ y + x \circ D(y)$ in every characteristic.
Setting $y = x$ in this condition gives $2D(x^2) = 2(D(x)x + xD(x))$.
So the two conditions agree when $2$ is invertible in $K$, and we use the $\circ$ form throughout.
When $\operatorname{char} K = 2$ and $A$ is commutative, $x \circ y = 0$, and every linear map satisfies the $\circ$ condition.

We need two more kinds.
A \emph{Lie derivation} satisfies $D([x,y]) = [D(x), y] + [x, D(y)]$ \citep[p.~183]{Martindale1964}, and we write $\LDer(A)$ for this space.
When $A$ is commutative, every bracket is zero, so every linear map is a Lie derivation.
For an automorphism $\alpha$ of $A$, an \emph{$\alpha$-derivation} satisfies $D(xy) = D(x)\alpha(y) + \alpha(x)D(y)$ \citep[\S1]{HosseiniKarizaki2019}.
We write $\ADer(A)$ for this space.

We now compare these spaces.
When $2$ is invertible, $xy = \tfrac12 (x \circ y) + \tfrac12 [x,y]$.
A derivation respects $xy$, so it respects both halves, and $\Der(A) \subseteq \JDer(A) \cap \LDer(A)$.
Both inclusions can be strict, as Examples~\ref{ex:dual-numbers} and~\ref{ex:matrix} show.

Section~\ref{sec:intro} recalled the results of Herstein, Cusack and Bre\v{s}ar on Jordan derivations, and the central term $t$ that Lie derivations must allow.
In the theorems of \citet[Theorem 2]{Martindale1964} and \citet[Theorem 5]{Bresar1993}, every Lie derivation has the form $D = d + t$.
Here $d$ is a derivation into the central closure.
Here the map $t$ is additive, takes values in the extended centroid, and vanishes on commutators.

We now recall these two notions for a prime ring $R$ \citep[p.~525]{Bresar1993}.
Take a nonzero ideal $U$ of $R$ and a map $f \colon U \to R$ of right $R$-modules.
We treat two such maps as equal if they agree on a nonzero ideal that lies in both domains.
We add and compose these maps, and this gives a ring $Q$ \citep[p.~525]{Bresar1993}.

The ring $R$ lies inside $Q$, since each element $a$ of $R$ acts as the map $x \mapsto ax$ \citep[p.~525]{Bresar1993}.
The centre $C$ of $Q$ is a field \citep[p.~525]{Bresar1993}.
We call $C$ the \emph{extended centroid} of $R$.
We call the ring $RC + C$ the \emph{central closure} of $R$.
Both notions go back to Martindale \citep[p.~525]{Bresar1993}.

We recall a polynomial way to unify these maps.
A \emph{multilinear polynomial} in variables $x_1, \dots, x_n$ that do not commute has the following form.
\[
f = \sum_{\varpi} \kappa_\varpi\, x_{\varpi(1)} \cdots x_{\varpi(n)}.
\]
Here $\varpi$ runs over the permutations of $\{1, \dots, n\}$, and the coefficients $\kappa_\varpi$ lie in $K$ \citep[\S1, (1.1)]{Benkovic2015}.

Fix a nonzero multilinear polynomial $f$ in $n$ variables.
An \emph{$f$-derivation} satisfies the identity
\[
D(f(x_1, \dots, x_n)) = \sum_{i=1}^n f(x_1, \dots, D(x_i), \dots, x_n)
\]
for all $x_1, \dots, x_n$ \citep[\S1, (1.2)]{Benkovic2015}.
Derivations, Jordan derivations and Lie derivations are the $f$-derivations for $f = x_1x_2$, $x_1 \circ x_2$ and $[x_1, x_2]$ \citep[\S1]{Benkovic2015}.
In each case $f$ uses only sums, products and scalars, so it makes sense in every algebra.

The twisted products of Section~\ref{sec:twisted} are different.
They have the form
\[
xy + \alpha(y)\alpha(x)
\]
for a fixed automorphism or anti-automorphism $\alpha$ of $A$.
Here the expression applies $\alpha$ to the variables, and $\alpha$ is not a polynomial operation.
So it is not of the form $f(x,y)$ for any polynomial $f$.
Hence the results on $f$-derivations do not apply directly.

Our examples include triangular algebras.
An \emph{$(A,B)$-bimodule} is a left $A$-module and right $B$-module $M$ with $(am)b = a(mb)$ for all $a \in A$, $m \in M$ and $b \in B$ \citep[Example 1.14]{Lam2001}.
We call it \emph{unital} if $1m = m = m1$ for all $m \in M$.

Let $R$ be a commutative ring, let $A$ and $B$ be unital $R$-algebras, and let $M$ be a unital $(A,B)$-bimodule.
The \emph{triangular algebra} $T = \mathrm{Tri}(A,M,B)$ is the set of matrices $\left(\begin{smallmatrix} a & m \\ 0 & b \end{smallmatrix}\right)$ with the matrix sum and product \citep[\S2]{Cheung2003}.
The module $M$ is \emph{faithful} as a left $A$-module if $aM = 0$ implies $a = 0$, and as a right $B$-module if $Mb = 0$ implies $b = 0$.
The idempotents $e_1 = \left(\begin{smallmatrix} 1 & 0 \\ 0 & 0 \end{smallmatrix}\right)$ and $e_2 = \left(\begin{smallmatrix} 0 & 0 \\ 0 & 1 \end{smallmatrix}\right)$ satisfy $e_2 T e_1 = 0$.
So $T$ is not prime when $M \neq 0$, and the prime-ring results above do not apply to triangular algebras.

We use three involutions of $A \otimes A$: the identity, the swap $\tau(x \otimes y) = y \otimes x$, and the $\alpha$-twisted map $\sigma_\alpha(x \otimes y) = \alpha(y) \otimes \alpha(x)$ of Section~\ref{sec:twisted}.

Let $\sigma \colon A \otimes A \to A \otimes A$ be a $K$-linear map with $\sigma^2 = \mathrm{id}$.
We write $x \cdot_\sigma y = xy + \mu(\sigma(x \otimes y))$.
We ask for no further axiom.
The derivations below are called braided-symmetric because $\sigma$ plays the role of the braiding operator in the work on braided derivations cited above.

\section{Braided-symmetric derivations}\label{sec:sigma-derivations}

This section defines braided-symmetric derivations and identifies them for the identity and the swap.

\begin{definition}\label{def:sigma-derivation}
Let $A$ be a $K$-algebra and let $\sigma$ be a $K$-linear involution of $A \otimes A$.
A \emph{braided-symmetric derivation} for $\sigma$ is a $K$-linear map $D \colon A \to A$ that satisfies identity~\eqref{eq:sigma-leibniz}, that is, $D(x \cdot_\sigma y) = D(x) \cdot_\sigma y + x \cdot_\sigma D(y)$ for all $x, y \in A$.
We write $\BSDer_\sigma(A)$ for the set of braided-symmetric derivations for $\sigma$.
This set is a $K$-subspace of $\mathrm{End}_K(A)$, since the condition is linear in $D$.
For an automorphism $\alpha$ of $A$, we write $\BSDer_\sigma(A)^\alpha$ for the subset of maps $D$ with $D\alpha = \alpha D$.
\end{definition}

The name is close to an older one, and the next remark separates the two.

\begin{remark}\label{rem:distinction-from-braided}
Braided derivations \citep[\S2]{Rozanski1996} \citep[\S4.1, Definition 14]{Huru2006} form a different notion from braided-symmetric derivations, although both use a braiding.
\citet[\S3, equation (3.9)]{Weber2020} defines a braided derivation $X$ of a braided commutative module algebra $A$ over a triangular Hopf algebra $(H, R)$ by a Leibniz rule for the product $ab$ itself:
\[
X(ab) = X(a)b + (R_1^{-1} \triangleright a)\bigl((R_2^{-1} \triangleright X)(b)\bigr).
\]
The braiding enters through the $R$-matrix.
Definition~\ref{def:sigma-derivation} asks instead for the ordinary Leibniz rule for the new product $\mu + \mu \circ \sigma$.

\end{remark}

Example~\ref{ex:not-associative} shows that the product $\cdot_\sigma$ can fail to be associative and unital, even when $A$ is commutative.

\begin{example}\label{ex:not-associative}
Let $A = K \times K$ with componentwise operations, and let $\alpha(a,b) = (b,a)$.
Then $\alpha$ is an automorphism with $\alpha^2 = \mathrm{id}$, and $\sigma_\alpha(x \otimes y) = \alpha(y) \otimes \alpha(x)$ is an involution of $A \otimes A$.
Since $A$ is commutative, $x \cdot_{\sigma_\alpha} y = xy + \alpha(xy)$.
For $x = (x_1, x_2)$ and $y = (y_1, y_2)$ this gives $x \cdot_{\sigma_\alpha} y = (x_1y_1 + x_2y_2)(1,1)$.
Put $p = (1,0)$ and $q = (0,1)$.
Then $(p \cdot_{\sigma_\alpha} p) \cdot_{\sigma_\alpha} q = (1,1) \cdot_{\sigma_\alpha} q = (1,1)$, while $p \cdot_{\sigma_\alpha} (p \cdot_{\sigma_\alpha} q) = p \cdot_{\sigma_\alpha} 0 = 0$.
The two bracketings are the two trees
\[
\begin{tikzpicture}[baseline=0.6cm,x=0.45cm,y=0.45cm,line width=0.5pt]
\node[font=\scriptsize] (a) at (0,0) {$p$}; \node[font=\scriptsize] (b) at (1,0) {$p$}; \node[font=\scriptsize] (c) at (2.4,0) {$q$};
\draw (a) -- (0.5,1.2) -- (b); \fill (0.5,1.2) circle (1.5pt);
\node[font=\scriptsize,left] at (0.4,1.3) {$(1,1)$};
\draw (0.5,1.2) -- (1.4,2.6) -- (c); \fill (1.4,2.6) circle (1.5pt); \draw (1.4,2.6) -- (1.4,3.3);
\node[font=\scriptsize,above] at (1.4,3.3) {$(1,1)$};
\end{tikzpicture}
\qquad\text{and}\qquad
\begin{tikzpicture}[baseline=0.6cm,x=0.45cm,y=0.45cm,line width=0.5pt]
\node[font=\scriptsize] (a) at (0,0) {$p$}; \node[font=\scriptsize] (b) at (1.4,0) {$p$}; \node[font=\scriptsize] (c) at (2.4,0) {$q$};
\draw (b) -- (1.9,1.2) -- (c); \fill (1.9,1.2) circle (1.5pt);
\node[font=\scriptsize,right] at (1.95,1.2) {$0$};
\draw (a) -- (1,2.6) -- (1.9,1.2); \fill (1,2.6) circle (1.5pt); \draw (1,2.6) -- (1,3.3);
\node[font=\scriptsize,above] at (1,3.3) {$0$};
\end{tikzpicture}
\]
where each dot is the product $\cdot_{\sigma_\alpha}$ and each label is the value on that edge.
So $\cdot_{\sigma_\alpha}$ is not associative.
Every product $u \cdot_{\sigma_\alpha} y$ lies in $K(1,1)$, so $u \cdot_{\sigma_\alpha} p \neq p$ for every $u$, and $\cdot_{\sigma_\alpha}$ has no unit.
\end{example}

The next remark records what Definition~\ref{def:sigma-derivation} says for the two involutions that use no structure of $A$.

\begin{remark}\label{rem:sigma-derivation-classification}
Let $A$ be a $K$-algebra and let $\sigma$ be an involution of $A \otimes A$.
Identity~\eqref{eq:sigma-leibniz} for $\sigma$ is the Leibniz rule for the product $\cdot_\sigma$.
So $\BSDer_\sigma(A)$ is the set of derivations of the nonassociative algebra $(A, \cdot_\sigma)$, in the sense of \citet[Part I]{Jacobson1937}.

For the identity, $x \cdot_{\mathrm{id}} y = 2xy$, so the identity reads
\[
2D(xy) = 2\bigl(D(x)y + xD(y)\bigr).
\]
If $\operatorname{char} K \neq 2$, we divide by $2$ and get $\BSDer_{\mathrm{id}}(A) = \Der(A)$.
If $\operatorname{char} K = 2$, then $x \cdot_{\mathrm{id}} y = 0$, so $\BSDer_{\mathrm{id}}(A)$ contains every linear map.
So it is larger than $\Der(A)$.
For example, take the linear map with $D(1) = 1$ and $D = 0$ on a complement of $K1$.
Then $D$ is not a derivation, since a derivation satisfies $D(1) = D(1 \cdot 1) = 2D(1) = 0$ in characteristic $2$.

For the swap $\tau$ we get $x \cdot_\tau y = xy + yx = x \circ y$.
So the defining identity is the $\circ$ form of the Jordan identity from Section~\ref{sec:preliminaries}.
Hence $\BSDer_\tau(A) = \JDer(A)$ in every characteristic.

Every derivation is a Jordan derivation \citep[p.~1003]{Bresar1988}.
So in every characteristic
\[
\Der(A) \subseteq \BSDer_\tau(A) = \JDer(A).
\]
When $\operatorname{char} K \neq 2$, this gives $\BSDer_{\mathrm{id}}(A) \subseteq \BSDer_\tau(A)$.
This is an inclusion between the spaces of two involutions on one algebra. % claim-ok: names the inclusion displayed just above
\end{remark}

The smallest nonreduced algebra shows both inclusions of Section~\ref{sec:preliminaries} at work.

\begin{example}[Dual numbers]\label{ex:dual-numbers}
Let $A = K[\varepsilon]/(\varepsilon^2)$, with basis $\{1, \varepsilon\}$.
A linear map $D$ is a derivation exactly when the Leibniz rule holds on the four pairs of basis elements.
Since $D(1) = 0$, the pairs $(1,\varepsilon)$ and $(\varepsilon,1)$ hold.
Write $D(\varepsilon) = \beta + c\varepsilon$ with $\beta, c \in K$.
The pair $(\varepsilon,\varepsilon)$ gives $0 = D(\varepsilon^2) = 2\varepsilon D(\varepsilon) = 2\beta\varepsilon$.
If $\operatorname{char} K \neq 2$, then $\beta = 0$, and every derivation has the form $D(a + b\varepsilon) = bc\varepsilon$, so $\dim_K \Der(A) = 1$.
If $\operatorname{char} K = 2$, then $\beta$ is free, and $\dim_K \Der(A) = 2$.
Since $A$ is commutative, every commutator is zero, so $\LDer(A) = \mathrm{End}_K(A)$, which has dimension $2 \cdot 2 = 4$.
Each derivation lies in $\BSDer_\tau(A)$ by Remark~\ref{rem:sigma-derivation-classification}.
If $\operatorname{char} K = 2$, then $x \circ y = 2xy = 0$ for all $x, y$, so $\JDer(A) = \mathrm{End}_K(A)$ has dimension $4$, which exceeds $\dim_K \Der(A) = 2$.
\end{example}

The remark above and the following theorem on triangular algebras give a corollary. % claim-ok: signpost to the corollary proved right after

\begin{theorem}[{\citep[Theorem 2.1]{ZhangYu2006}}]\label{thm:zhang-yu}
Let $R$ be a $2$-torsion free commutative ring, let $A$ and $B$ be unital $R$-algebras, and let $M$ be a unital $(A,B)$-bimodule that is faithful as a left $A$-module and as a right $B$-module. Then every Jordan derivation $D \colon \mathrm{Tri}(A,M,B) \to \mathrm{Tri}(A,M,B)$ is a derivation.
\end{theorem}

In the next corollary, maps on $T = \mathrm{Tri}(A,M,B)$ are $R$-linear and the braided-symmetric identity is read over $R$.
It follows from Remark~\ref{rem:sigma-derivation-classification} and Theorem~\ref{thm:zhang-yu}.

\begin{corollary}\label{cor:triangular-jordan}
Let $R$ be a $2$-torsion free commutative ring, let $A$ and $B$ be unital $R$-algebras, and let $M$ be a unital $(A,B)$-bimodule that is faithful as a left $A$-module and as a right $B$-module.
Suppose also that $\mathrm{Tri}(A,M,B)$ is $2$-torsion free.
Then $\BSDer_\tau(\mathrm{Tri}(A,M,B)) = \Der(\mathrm{Tri}(A,M,B))$.
\end{corollary}

Proposition~\ref{prop:base-change} below treats extension of the base field.
For a field extension $L$ of $K$ and a $K$-vector space $V$, write $V_L = L \otimes_K V$.
This is an $L$-vector space. % claim-ok: the tensor product of a $K$-vector space with $L$ over $K$ is an $L$-vector space by construction
For a $K$-algebra $A$, the space $A_L$ is an $L$-algebra with $(c \otimes x)(c' \otimes y) = cc' \otimes xy$ \citep[Chapter 2, \S5, p.~74]{Lam2001}.

We now define the involution $\sigma_L$ of $A_L \otimes_L A_L$.
It is an $L$-linear map of $A_L \otimes_L A_L$ into itself with $\sigma_L^2 = \mathrm{id}$.

Define
\[
\Psi \colon A_L \otimes_L A_L \to L \otimes_K (A \otimes_K A), \qquad
(c \otimes x) \otimes_L (c' \otimes y) \longmapsto cc' \otimes (x \otimes y).
\]
Since this assignment is $L$-bilinear, it defines an $L$-linear map.
It is an isomorphism, because its inverse sends $c \otimes (x \otimes y)$ to $(c \otimes x) \otimes_L (1 \otimes y)$.

The map $\mathrm{id}_L \otimes \sigma$ is an $L$-linear map of $L \otimes_K (A \otimes_K A)$ into itself, and $(\mathrm{id}_L \otimes \sigma)^2 = \mathrm{id}_L \otimes \sigma^2 = \mathrm{id}$.
We define $\sigma_L$ as the conjugate of $\mathrm{id}_L \otimes \sigma$ by $\Psi$:
\[
\sigma_L = \Psi^{-1} \circ (\mathrm{id}_L \otimes \sigma) \circ \Psi.
\]
For $c, c' \in L$ and $x, y \in A$ it reads
\[
\sigma_L\bigl((c \otimes x) \otimes_L (c' \otimes y)\bigr) = \Psi^{-1}\bigl(cc' \otimes \sigma(x \otimes y)\bigr).
\]

Since $\Psi \Psi^{-1} = \mathrm{id}$, we get
\[
\sigma_L^2 = \Psi^{-1} (\mathrm{id}_L \otimes \sigma)(\mathrm{id}_L \otimes \sigma) \Psi = \Psi^{-1} (\mathrm{id}_L \otimes \sigma^2) \Psi = \mathrm{id}.
\]
So $\sigma_L$ is an involution of $A_L \otimes_L A_L$.

The product $\mu_L$ of $A_L$ satisfies $\mu_L = (\mathrm{id}_L \otimes \mu) \circ \Psi$.
So $\Psi$ carries $\mu_L$ and $\sigma_L$ to $\mathrm{id}_L \otimes \mu$ and $\mathrm{id}_L \otimes \sigma$.
Hence $(1 \otimes x) \cdot_{\sigma_L} (1 \otimes y) = 1 \otimes (x \cdot_\sigma y)$ for $x, y \in A$.
Thus the product $\cdot_{\sigma_L}$ on $A_L$ extends the product $\cdot_\sigma$ on $A$.

The proof uses three results from the literature, which we recall. % claim-ok: signpost for the three cited theorems below
For a $K$-vector space $E$ we write $E_L = L \otimes_K E$, and for a $K$-linear map $f$ we write $f_L = \mathrm{id}_L \otimes f$.

\begin{theorem}[{\cite[Chapter II, \S7, no.~9, Proposition 18 and its Corollary]{Bourbaki1989}}]\label{thm:ext-exact}
Let $L$ be an extension field of the field $K$. If $E' \to E \to E''$ is an exact sequence of $K$-vector spaces and $K$-linear maps, then $E'_L \to E_L \to E''_L$ is exact.
Hence, for every $K$-linear map $f \colon E \to F$, we have $\ker(f_L) = (\ker f)_L$ and $\operatorname{im}(f_L) = (\operatorname{im} f)_L$, up to canonical isomorphism.
\end{theorem}

\begin{theorem}[{\cite[Theorem 3.1]{Rotman2009}}]\label{thm:free-projective}
Every free module over a ring is projective.
\end{theorem}

\begin{theorem}[{\cite[Chapter II, \S5, no.~3, Proposition 7(ii)]{Bourbaki1989}}]\label{thm:hom-base-change}
Let $A$ be a commutative ring, let $B$ be a ring, let $\rho \colon A \to B$ be a ring homomorphism with image in the centre of $B$, and let $E$, $F$ be $A$-modules.
If $E$ is a finitely generated projective $A$-module, then the canonical homomorphism $B \otimes_A \mathrm{Hom}_A(E, F) \to \mathrm{Hom}_B(B \otimes_A E,\, B \otimes_A F)$, given by $b \otimes u \mapsto b\,(\mathrm{id}_B \otimes u)$, is bijective.
\end{theorem}

If $V$ is a finite-dimensional $K$-vector space, then $V$ is free with a finite basis, so it is a finitely generated projective $K$-module by Theorem~\ref{thm:free-projective}.
So Theorem~\ref{thm:hom-base-change}, with $A = K$ and $B = L$, gives a bijection $L \otimes_K \mathrm{Hom}_K(V, W) \to \mathrm{Hom}_L(V_L, W_L)$ for every $K$-vector space $W$.

For a $K$-algebra $A$ with involution $\sigma$ of $A \otimes A$, let $m_\sigma \colon A \otimes A \to A$ be the product $x \otimes y \mapsto x \cdot_\sigma y$.
Define the $K$-linear map
\[
\Phi \colon \mathrm{End}_K(A) \to \mathrm{Hom}_K(A \otimes A, A), \qquad
\Phi(D) = D \circ m_\sigma - m_\sigma \circ (D \otimes \mathrm{id} + \mathrm{id} \otimes D).
\]
So $\Phi(D)$ is the difference of the two sides of identity~\eqref{eq:sigma-leibniz}, and $\BSDer_\sigma(A) = \ker \Phi$.
Write $\Phi_L$ for the same map for $A_L$ and $\sigma_L$, so $\BSDer_{\sigma_L}(A_L) = \ker \Phi_L$.

\begin{proposition}\label{prop:base-change}
Let $K$ be a field, let $A$ be a finite-dimensional $K$-algebra, and let $\sigma$ be an involution of $A \otimes A$.
Let $L$ be a field extension of $K$.
For $D \in \mathrm{End}_K(A)$ and $c \in L$, let $c \otimes D$ act on $A_L$ by $(c \otimes D)(c' \otimes x) = cc' \otimes D(x)$ for $c' \in L$ and $x \in A$.
This rule defines isomorphisms of $L$-vector spaces
\[
\beta \colon L \otimes_K \mathrm{End}_K(A) \to \mathrm{End}_L(A_L), \qquad
\gamma \colon L \otimes_K \mathrm{Hom}_K(A \otimes A, A) \to \mathrm{Hom}_L(A_L \otimes_L A_L, A_L)
\]
with $\Phi_L \circ \beta = \gamma \circ (\mathrm{id}_L \otimes \Phi)$.
So $\beta$ and $\gamma$ carry the exact sequence
\[
0 \to L \otimes_K \BSDer_\sigma(A) \to L \otimes_K \mathrm{End}_K(A) \xrightarrow{\ \mathrm{id}_L \otimes \Phi\ } L \otimes_K \mathrm{Hom}_K(A \otimes A, A)
\]
isomorphically onto the exact sequence
\[
0 \to \BSDer_{\sigma_L}(A_L) \to \mathrm{End}_L(A_L) \xrightarrow{\ \Phi_L\ } \mathrm{Hom}_L(A_L \otimes_L A_L, A_L).
\]
In particular $L \otimes_K \BSDer_\sigma(A) \cong \BSDer_{\sigma_L}(A_L)$ and $\dim_L \BSDer_{\sigma_L}(A_L) = \dim_K \BSDer_\sigma(A)$.
\end{proposition}

\begin{proof}
First, the spaces $A$ and $A \otimes_K A$ are finite-dimensional.
So $\beta$ is an isomorphism by the bijection above with $V = W = A$.
Also $\gamma$ is an isomorphism by the bijection above with $V = A \otimes_K A$ and $W = A$, since $(A \otimes_K A)_L = A_L \otimes_L A_L$.

Second, we check $\Phi_L \circ \beta = \gamma \circ (\mathrm{id}_L \otimes \Phi)$.
Both sides are $L$-linear.
So it is enough to evaluate them on $c \otimes D$ and on $(c' \otimes x) \otimes_L (c'' \otimes y)$, for $c, c', c'' \in L$, $D \in \mathrm{End}_K(A)$ and $x, y \in A$.
By the construction of $\sigma_L$, we have $m_{\sigma_L}\bigl((c' \otimes x) \otimes_L (c'' \otimes y)\bigr) = c'c'' \otimes m_\sigma(x \otimes y)$.
Hence both sides give
\[
cc'c'' \otimes \Bigl(D\bigl(m_\sigma(x \otimes y)\bigr) - m_\sigma\bigl(D(x) \otimes y + x \otimes D(y)\bigr)\Bigr) = cc'c'' \otimes \Phi(D)(x \otimes y).
\]

Third, $\ker \Phi = \BSDer_\sigma(A)$.
By Theorem~\ref{thm:ext-exact}, $\ker(\mathrm{id}_L \otimes \Phi) = L \otimes_K \BSDer_\sigma(A)$.
By the second step, $\beta$ and $\gamma$ turn $\mathrm{id}_L \otimes \Phi$ into $\Phi_L$.
So $\beta$ maps $L \otimes_K \BSDer_\sigma(A)$ onto $\ker \Phi_L = \BSDer_{\sigma_L}(A_L)$.
\end{proof}


\begin{remark}\label{rem:finite-needed}
The finite-dimensional hypothesis in Proposition~\ref{prop:base-change} cannot be dropped.
Take $\sigma = -\mathrm{id}$.
Then $\sigma^2 = \mathrm{id}$, so $\sigma$ is an involution of $A \otimes A$, and $x \cdot_\sigma y = 0$ for all $x, y$.
So every $K$-linear map is a braided-symmetric derivation, both for $A$ and for $A_L$.

Let $A = K[x]$ and $L = K(t)$, where $t$ is transcendental over $K$.
Define $D' \in \mathrm{End}_L(A_L)$ by $D'(x^i) = t^i x^i$ for all $i \geq 0$.
Suppose $D' = \sum_{r=1}^s c_r \otimes D_r$ for some $s \geq 1$, with $c_r \in L$ and $D_r \in \mathrm{End}_K(A)$.
Since $D_r$ is $K$-linear and $A = K[x]$, the coefficient of $x^i$ in $D_r(x^i)$ lies in $K$.
So $t^i$ lies in the $K$-span of $c_1, \dots, c_s$, for every $i$.
Since $t$ is transcendental over $K$, the powers $1, t, t^2, \dots$ are linearly independent over $K$, so no finite span contains them all.
Hence $D'$ is not in the image, and the map of Proposition~\ref{prop:base-change} is not onto.
\end{remark}

The fixed part $\BSDer_\sigma(A)^\alpha$ is also a kernel, for an automorphism $\alpha$ of $A$. % claim-ok: by definition of the fixed part, shown in the display below
It is the kernel of the $K$-linear map
\[
\Phi_\alpha \colon \mathrm{End}_K(A) \to \mathrm{Hom}_K(A \otimes A, A) \oplus \mathrm{End}_K(A), \qquad D \mapsto \bigl(\Phi(D),\, D\alpha - \alpha D\bigr).
\]
Over $L$, the same map is built from $\alpha_L = \mathrm{id}_L \otimes \alpha$.
The isomorphisms $\beta$ and $\gamma$ identify it with $\mathrm{id}_L \otimes \Phi_\alpha$. % claim-ok: the check of the proposition applied to the second component, where alpha_L corresponds to alpha under beta
By Theorem~\ref{thm:ext-exact}, kernels commute with $L \otimes_K -$.
Hence
\[
\dim_L \BSDer_{\sigma_L}(A_L)^{\alpha_L} = \dim_K \BSDer_\sigma(A)^\alpha.
\]

Here is an example.
Let $K = \mathbb{R}$, $A = \mathbb{C}$, and let $\alpha$ be complex conjugation.
Then $A$ is commutative and reduced, so Theorem~\ref{thm:twisted-commutative} below gives
\[
\BSDer_{\sigma_\alpha}(A)^\alpha = \Der(A)^\alpha.
\]

Every $\mathbb{R}$-linear derivation $D$ of $\mathbb{C}$ is zero, since $2iD(i) = D(i^2) = D(-1) = 0$.
Now extend to $L = \mathbb{C}$.
The algebra $A_L = \mathbb{C} \otimes_\mathbb{R} \mathbb{C}$ is isomorphic to $\mathbb{C} \times \mathbb{C}$, and $\alpha_L$ becomes the swap of the two factors.
Both spaces are zero, as Proposition~\ref{prop:base-change} predicts.

\section{Twisted involutions and the Leibniz defect}\label{sec:twisted}

Remark~\ref{rem:sigma-derivation-classification} treated the identity and the swap, which use no structure of $A$.
This section twists the swap by an involution $\alpha$ of $A$ and treats it with one tool, the Leibniz defect.

We recall the notions used here.
An \emph{automorphism} of $A$ is a bijective $K$-linear map $\alpha \colon A \to A$ with $\alpha(xy) = \alpha(x)\alpha(y)$ and $\alpha(1) = 1$.
An \emph{anti-automorphism} of $A$ is a bijective $K$-linear map $\alpha \colon A \to A$ with $\alpha(xy) = \alpha(y)\alpha(x)$ and $\alpha(1) = 1$.
By $(xy)^{T} = y^{T}x^{T}$, the transpose is an anti-automorphism of $M_n(K)$.
In both cases $\alpha$ is \emph{involutive} if $\alpha^2 = \mathrm{id}$.
An involutive anti-automorphism is an involution of the ring $A$ in the sense of \citet[\S1.2]{Herstein1976}, that is, a map $*$ with $a^{**} = a$, $(a+b)^* = a^* + b^*$ and $(ab)^* = b^*a^*$.

We need $\alpha^2 = \mathrm{id}$ for the following reason.
Applying $\sigma_\alpha$ below twice gives
\[
x \otimes y \mapsto \alpha^2(x) \otimes \alpha^2(y),
\]
so $\sigma_\alpha$ is an involution of $A \otimes A$ exactly when $\alpha$ is involutive.
When $2$ is invertible, an involutive $\alpha$ splits $A = A_+ \oplus A_-$, where $A_\pm = \{x \in A : \alpha(x) = \pm x\}$.
Indeed, for $x \in A$,
\[
x = \tfrac12\bigl(x + \alpha(x)\bigr) + \tfrac12\bigl(x - \alpha(x)\bigr).
\]
We call the elements of $A_+$ \emph{symmetric} and the elements of $A_-$ \emph{skew}.

Let $\alpha$ be involutive, of either kind.
Then $\sigma_\alpha(x \otimes y) = \alpha(y) \otimes \alpha(x)$ is a $K$-linear involution of $A \otimes A$.
So $\sigma_\alpha$ defines a symmetrized product as in Section~\ref{sec:preliminaries}.
The product it defines is
\begin{equation}\label{eq:twisted-product}
x \cdot_{\sigma_\alpha} y = xy + \alpha(y)\alpha(x).
\end{equation}
For $\alpha = \mathrm{id}$ the product $\mu_{\sigma_\alpha}$ is the Jordan product.
If $\alpha$ is an automorphism, then $\alpha(y)\alpha(x) = \alpha(yx)$.
If $\alpha$ is an anti-automorphism, then $\alpha(y)\alpha(x) = \alpha(xy)$.

\begin{definition}\label{def:defect}
The \emph{Leibniz defect} of a linear map $D \colon A \to A$ is the bilinear map $E_D \colon A \times A \to A$ given, for $x, y \in A$, by
\[
E_D(x,y) = D(xy) - D(x)y - xD(y).
\]
\end{definition}

The name is ours, and the map is classical: \citet[\S2]{Cusack1975} writes it $d(a,b) = D(ab) - aDb - (Da)b$ and studies Jordan derivations through it.
We recall the Hochschild complex of $A$ with coefficients in $A$ \citep[\S9.1]{Weibel1994}.
An \emph{$n$-cochain} is a $K$-multilinear map from $A^n$ to $A$.
So a $0$-cochain is an element of $A$, and a $1$-cochain is a linear map $D \colon A \to A$.
The \emph{coboundary} $\delta$ sends each $n$-cochain to an $(n+1)$-cochain, and $\delta \circ \delta = 0$.
An \emph{$n$-cocycle} is an $n$-cochain $f$ with $\delta f = 0$.
The $n$-th \emph{Hochschild cohomology} $HH^n(A, A)$ is the space of $n$-cocycles modulo the images $\delta g$ of the $(n-1)$-cochains $g$ \citep[\S9.1]{Weibel1994}.

On a $1$-cochain $D$ the coboundary is $(\delta D)(x,y) = xD(y) - D(xy) + D(x)y$ \citep[\S2, Definition 2.1]{Hochschild1945} \citep[\S9.1]{Weibel1994}.
So $E_D = -\delta D$.
Thus $D$ is a derivation exactly when $E_D = 0$, so derivations are the $1$-cocycles.
Since $\delta \circ \delta = 0$, the defect is a $2$-cocycle: for all $a, b, c \in A$,
\begin{equation}\label{eq:cocycle}
aE_D(b,c) - E_D(ab,c) + E_D(a,bc) - E_D(a,b)c = 0.
\end{equation}
This says that $E_D$ is a cocycle of the Hochschild complex \citep[\S9.1]{Weibel1994}.

The defect also describes the other classes of Section~\ref{sec:preliminaries}.
Expanding $D(x \circ y)$ shows that $D$ is a Jordan derivation exactly when $E_D(x,y) + E_D(y,x) = 0$ for all $x, y$.
In the same way, $D$ is a Lie derivation exactly when $E_D(x,y) = E_D(y,x)$ for all $x, y$.

\begin{proposition}\label{prop:defect}
Let $A$ be a unital associative $K$-algebra, let $\alpha$ be an involutive automorphism or anti-automorphism of $A$, and let $D \colon A \to A$ be a linear map with $D\alpha = \alpha D$.
\begin{enumerate}[label=(\roman*)]
\item If $\alpha$ is an automorphism, then $D \in \BSDer_{\sigma_\alpha}(A)$ if and only if $E_D(x,y) + \alpha(E_D(y,x)) = 0$ for all $x, y \in A$.
\item If $\alpha$ is an anti-automorphism, then $D \in \BSDer_{\sigma_\alpha}(A)$ if and only if $E_D(x,y) + \alpha(E_D(x,y)) = 0$ for all $x, y \in A$.
\end{enumerate}
In both cases every derivation that commutes with $\alpha$ lies in $\BSDer_{\sigma_\alpha}(A)^\alpha$.
\end{proposition}

\begin{proof}
Write $E = E_D$.
By Remark~\ref{rem:sigma-derivation-classification}, $D \in \BSDer_{\sigma_\alpha}(A)$ means $D(x \cdot_{\sigma_\alpha} y) = D(x) \cdot_{\sigma_\alpha} y + x \cdot_{\sigma_\alpha} D(y)$ for all $x, y$.
We compute the difference of the two sides in each case.

For (i), the product~\eqref{eq:twisted-product} is $x \cdot_{\sigma_\alpha} y = xy + \alpha(yx)$.
Since $D\alpha = \alpha D$, the left side is $D(xy) + \alpha(D(yx))$.
The right side is $D(x)y + \alpha(yD(x)) + xD(y) + \alpha(D(y)x)$.
We group the terms without $\alpha$ and the terms with $\alpha$.
The first group gives $D(xy) - D(x)y - xD(y) = E(x,y)$.
Since $\alpha$ is linear, the second group gives $\alpha\bigl(D(yx) - D(y)x - yD(x)\bigr) = \alpha(E(y,x))$.
So the left side minus the right side is $E(x,y) + \alpha(E(y,x))$, which proves (i).

The proof of (ii) follows the same steps as the proof of (i).
Here the product~\eqref{eq:twisted-product} is $x \cdot_{\sigma_\alpha} y = xy + \alpha(xy)$.
Since $D\alpha = \alpha D$, the left side is
\[
D(xy) + \alpha(D(xy)),
\]
and the right side is
\[
D(x)y + \alpha(D(x)y) + xD(y) + \alpha(xD(y)).
\]
The same grouping gives the difference
\[
E(x,y) + \alpha(E(x,y)),
\]
which proves (ii).

If $D$ is a derivation, then $E = 0$, and both conditions hold.
\end{proof}

We now read the two conditions, assuming that $2$ is invertible.

In case (ii), the equation $E_D + \alpha(E_D) = 0$ says $\alpha(E_D) = -E_D$.
So $E_D$ takes values in $A_-$, the skew elements of $\alpha$.
In case (i), we write $E_D = E_{\mathrm{s}} + E_{\mathrm{a}}$.
Here $E_{\mathrm{s}}(x,y) = \tfrac12(E_D(x,y) + E_D(y,x))$ is symmetric.
And $E_{\mathrm{a}}(x,y) = \tfrac12(E_D(x,y) - E_D(y,x))$ is antisymmetric.

We swap $x$ and $y$ in the condition of (i).
Then we take the sum and the difference of the two equations.
So the condition holds exactly when $E_{\mathrm{s}}$ takes values in $A_-$ and $E_{\mathrm{a}}$ takes values in $A_+$.
Now let $\alpha = \mathrm{id}$.
Then $A_- = 0$ and $A_+ = A$, so the condition says $E_{\mathrm{s}} = 0$.
This is the Jordan condition, so Proposition~\ref{prop:defect}(i) recovers the swap case of Remark~\ref{rem:sigma-derivation-classification}.

\begin{remark}\label{rem:inner-twisted}
Let $\alpha$ be an involutive automorphism of $A$, and let $D \in \Der(A)$.
Then $E_D = 0$.
By~\eqref{eq:twisted-product}, the two sides of the $\sigma_\alpha$-identity differ by
\[
D(\alpha(yx)) - \alpha(D(yx)).
\]
Since $A$ is unital, the products $yx$ span $A$.
So $D \in \BSDer_{\sigma_\alpha}(A)$ if and only if $D\alpha = \alpha D$, and then $D \in \BSDer_{\sigma_\alpha}(A)^\alpha$.

Next we look at the inner derivations $\operatorname{ad}_g$ in this class. % claim-ok: signpost
The condition $\operatorname{ad}_g \alpha = \alpha \operatorname{ad}_g$ means
\[
[g - \alpha(g), \alpha(x)] = 0 \qquad \text{for all } x.
\]
Since $\alpha$ is bijective, this holds if and only if $g - \alpha(g) \in Z(A)$.
Let $\Inn_\alpha(A)$ be the space of the maps $\operatorname{ad}_g$ with $g - \alpha(g) \in Z(A)$.
These are the inner derivations in $\BSDer_{\sigma_\alpha}(A)$, so $\Inn_\alpha(A) = \Inn(A) \cap \BSDer_{\sigma_\alpha}(A)^\alpha$.
The quotient $\BSDer_{\sigma_\alpha}(A)^\alpha / \Inn_\alpha(A)$ is the space of \emph{twisted outer derivations}.
For $\alpha = \mathrm{id}$ the quotient is $\Der(A)/\Inn(A) = HH^1(A,A)$, by Lemma~\ref{lem:weibel}.
\end{remark}

\citet[Corollary 5]{Cusack1975} and \citet[Theorem 1]{Bresar1988} proved that a Jordan derivation of a $2$-torsion free semiprime ring is a derivation.

\begin{theorem}[{\citep[Theorem 1]{Bresar1988}}]\label{thm:bresar}
Let $R$ be a $2$-torsion free semiprime ring. Then every Jordan derivation $D \colon R \to R$ is a derivation.
\end{theorem}

\begin{corollary}\label{cor:twisted-semiprime}
Let $A$ be a $2$-torsion-free semiprime $K$-algebra with an involutive automorphism $\alpha$.
Then $\BSDer_{\sigma_\alpha}(A)^\alpha \cap \JDer(A) = \Der(A)^\alpha$, where $\Der(A)^\alpha$ is the set of derivations commuting with $\alpha$.
\end{corollary}

\begin{proof}
Every derivation is a Jordan derivation by Remark~\ref{rem:sigma-derivation-classification}.
With Proposition~\ref{prop:defect}, this gives $\Der(A)^\alpha \subseteq \BSDer_{\sigma_\alpha}(A)^\alpha \cap \JDer(A)$.
Conversely, every Jordan derivation of $A$ is a derivation by Theorem~\ref{thm:bresar}, which applies since a $K$-linear map is additive.
So a map in the intersection is a derivation that commutes with $\alpha$.
\end{proof}

\section{Reduced commutative algebras and twisted derivations}\label{sec:twisted-commutative}

For commutative algebras the Jordan hypothesis in Corollary~\ref{cor:twisted-semiprime} can be dropped.
Recall from Section~\ref{sec:preliminaries} that a reduced commutative algebra has no nonzero nilpotent element.

\begin{theorem}\label{thm:twisted-commutative}
Let $A$ be a commutative reduced $K$-algebra with $2$ invertible in $K$, and let $\alpha$ be an involutive automorphism of $A$.
Then $\BSDer_{\sigma_\alpha}(A)^\alpha = \Der(A)^\alpha$.
\end{theorem}

\begin{proof}
Every derivation commuting with $\alpha$ lies in $\BSDer_{\sigma_\alpha}(A)^\alpha$ by Proposition~\ref{prop:defect}.
Conversely, let $D \in \BSDer_{\sigma_\alpha}(A)^\alpha$.
Write $E = E_D$.
Since $A$ is commutative, $E$ is symmetric.
So $E$ takes values in $A_-$ by Proposition~\ref{prop:defect}(i).

Since $D$ commutes with $\alpha$, it maps each of $A_+$ and $A_-$ into itself.
So if $x$ and $y$ lie in the same part, then $E(x,y)$ lies in
\[
A_+ \cap A_- = 0.
\]

Now let $x \in A_+$ and $y \in A_-$, and put $z = E(x,y)$.
Then $z$ lies in $A_-$, and so do $xy$, $zx$ and $xz$, while $zy$ lies in $A_+$.
Apply the cocycle identity~\eqref{eq:cocycle} to $(a,b,c) = (z,x,y)$:
\[
zE(x,y) - E(zx,y) + E(z,xy) - E(z,x)y = 0.
\]
The pairs $(zx,y)$ and $(z,xy)$ lie in $A_-$, so $E$ vanishes on them, and we get $z^2 = E(z,x)y$.
Apply~\eqref{eq:cocycle} to $(a,b,c) = (x,z,y)$:
\[
xE(z,y) - E(xz,y) + E(x,zy) - E(x,z)y = 0.
\]
The pairs $(z,y)$ and $(xz,y)$ lie in $A_-$, and the pair $(x,zy)$ lies in $A_+$, so $E$ vanishes on all three and $E(x,z)y = 0$.
Since $E$ is symmetric, $z^2 = E(z,x)y = E(x,z)y = 0$.
So $z = 0$ because $A$ is reduced.
Since $E$ is bilinear and symmetric and $A = A_+ \oplus A_-$, we get $E = 0$.
Hence $D$ is a derivation.
\end{proof}

The mixed-parity step uses the cocycle identity twice. % claim-ok: shown in the proof above
Example~\ref{ex:twisted-nilpotent} shows that the reduced hypothesis cannot be dropped, and Example~\ref{ex:twisted-char-two} shows the same for the invertibility of $2$.

\begin{example}\label{ex:twisted-nilpotent}
Let $A = \mathbb{F}_5[u,v]/(u^2, v^2)$.
It has basis $\{1, u, v, uv\}$.
Let $\alpha$ fix $u$ and send $v$ to $-v$.
Let $D$ send $uv$ to $uv$ and the other three basis elements to $0$.
Then $D$ commutes with $\alpha$.
It satisfies the Leibniz rule for $x \cdot_{\sigma_\alpha} y$ on all basis pairs.
So $D \in \BSDer_{\sigma_\alpha}(A)^\alpha$.
But $D$ is not a derivation, since
\[
D(uv) = uv \qquad \text{while} \qquad D(u)v + uD(v) = 0.
\]
Direct computation on the four basis elements verifies all three statements. % claim-ok: each product of two basis elements of F_5[u,v]/(u^2,v^2) is read off from the multiplication table
Here $z = E(u,v) = uv$ is a nonzero element with $z^2 = 0$.
\end{example}


\begin{example}\label{ex:twisted-char-two}
Let $A = \mathbb{F}_2 \times \mathbb{F}_2$, which is commutative and reduced, and let $\alpha$ swap the two factors.
Put $p = (1,0)$, $q = (0,1)$ and $e = (1,1)$, and let $D$ be the linear map with $D(p) = D(q) = e$, so that $D(e) = 0$.
Then $D$ commutes with $\alpha$.
On the basis pairs $(p,p)$, $(q,q)$ and $(p,q)$, both sides of the Leibniz rule for $x \cdot_{\sigma_\alpha} y = xy + \alpha(xy)$ are zero.
So $D \in \BSDer_{\sigma_\alpha}(A)^\alpha$.
But $D$ is not a derivation: $D(p) = D(p^2) = 2pD(p) = 0$ would hold for a derivation, while $D(p) = e$.
\end{example}

Our second natural candidate for the twisted class is the space $\ADer(A)$ of $\alpha$-derivations.
It was defined for automorphisms in Section~\ref{sec:preliminaries}.
It differs from $\BSDer_{\sigma_\alpha}(A)^\alpha$ in both directions, and one algebra shows both failures.

\begin{proposition}\label{prop:twisted}
Let $A = \mathbb{F}_3[\varepsilon]/(\varepsilon^3)$, and let $\alpha$ be the automorphism with $\alpha(\varepsilon) = -\varepsilon$.
Then
\[
\ADer(A) \not\subseteq \BSDer_{\sigma_\alpha}(A)^\alpha
\qquad\text{and}\qquad
\BSDer_{\sigma_\alpha}(A)^\alpha \not\subseteq \ADer(A).
\]
\end{proposition}

\begin{proof}
The algebra $A$ has basis $\{1, \varepsilon, \varepsilon^2\}$.
The map $\alpha$ fixes $1$ and $\varepsilon^2$ and sends $\varepsilon$ to $-\varepsilon$.
So it is an involutive automorphism with $A_+ = \mathbb{F}_3 1 + \mathbb{F}_3\varepsilon^2$ and $A_- = \mathbb{F}_3\varepsilon$.
Since $A$ is commutative, $E_D$ is symmetric for every linear map $D$.
Let $D$ be a linear map with $D\alpha = \alpha D$.
By Proposition~\ref{prop:defect}(i), $D$ lies in $\BSDer_{\sigma_\alpha}(A)$ exactly when $E_D$ takes values in $A_-$.
We give one map for each non-inclusion.
Both maps are diagonal on the basis, so both commute with $\alpha$.

For the first non-inclusion, let $D(a + b\varepsilon + c\varepsilon^2) = b\varepsilon + c\varepsilon^2$.
We check that $D$ is an $\alpha$-derivation on pairs of basis elements.
Pairs that contain $1$ give $0$ on both sides, since $D(1) = 0$ and $\alpha(1) = 1$.
For the pair $(\varepsilon, \varepsilon)$, the left side is $D(\varepsilon^2) = \varepsilon^2$.
The right side is $D(\varepsilon)\alpha(\varepsilon) + \alpha(\varepsilon)D(\varepsilon) = -2\varepsilon^2 = \varepsilon^2$.
The remaining pairs have product $0$, and each term on the right side is a multiple of $\varepsilon^3 = 0$.
So $D \in \ADer(A)$.
On the other hand, $E_D(\varepsilon, \varepsilon) = D(\varepsilon^2) - 2\varepsilon D(\varepsilon) = \varepsilon^2 - 2\varepsilon^2 = 2\varepsilon^2$.
This element is nonzero and lies in $A_+$, so it does not lie in $A_-$.
Hence $D \notin \BSDer_{\sigma_\alpha}(A)^\alpha$.

The second non-inclusion uses a derivation.

Let $d$ be the linear map
\[
d(a + b\varepsilon + c\varepsilon^2) = b\varepsilon + 2c\varepsilon^2.
\]
It is a derivation, since $d(\varepsilon^i) = i\varepsilon^i$ and exponents add under multiplication.
So $E_d = 0$, and $d \in \BSDer_{\sigma_\alpha}(A)^\alpha$.
On the other hand,
\[
d(\varepsilon^2) - d(\varepsilon)\alpha(\varepsilon) - \alpha(\varepsilon)d(\varepsilon) = 2\varepsilon^2 + 2\varepsilon^2 = \varepsilon^2 \neq 0.
\]
Hence $d \notin \ADer(A)$.
\end{proof}

The two classes can also meet, and the next example shows both cases.

\begin{example}\label{ex:twisted-meet}
Let $A = M_2(\mathbb{F}_3)$ and let $\alpha$ be the transpose, an anti-automorphism.
Then $\ADer(A)$ is not defined, since it needs an automorphism.
Take an inner derivation $\operatorname{ad}_g$ with $g$ skew for the transpose.
It commutes with $\alpha$ and has $E = 0$, so it lies in $\BSDer_{\sigma_\alpha}(A)^\alpha$ by Proposition~\ref{prop:defect}.

Now let $A = K[\varepsilon]/(\varepsilon^2)$ with $\alpha(\varepsilon) = -\varepsilon$.
The map $D(a + b\varepsilon) = b\varepsilon$ is a derivation and an $\alpha$-derivation, and it commutes with $\alpha$.
So it lies in both $\ADer(A)$ and $\BSDer_{\sigma_\alpha}(A)^\alpha$ by Proposition~\ref{prop:defect}.
\end{example}

\section{The antisymmetrized product and Lie derivations}\label{sec:lie}

This section replaces the sum $\mu + \mu \circ \sigma$ by the difference $\mu - \mu \circ \sigma$.
For the swap the difference is the commutator, and its derivations are the Lie derivations.

\begin{definition}\label{def:anti-sigma-derivation}
The \emph{$\sigma$-antisymmetrized product} on a $K$-algebra $A$ is $x \cdot_\sigma^- y = \mu(x \otimes y) - \mu(\sigma(x \otimes y))$, where $\sigma$ is an involution of $A \otimes A$.
A \emph{braided-antisymmetric derivation} is a $K$-linear map $D \colon A \to A$ with $D(x \cdot_\sigma^- y) = D(x) \cdot_\sigma^- y + x \cdot_\sigma^- D(y)$.
\end{definition}

We now identify these derivations for the swap.
It also gives the easy half of the theorems of Martindale and Bre\v{s}ar.

\begin{proposition}\label{prop:lie-classification}
Let $A$ be an associative $K$-algebra.
\begin{enumerate}[label=(\roman*)]
\item The braided-antisymmetric derivations for $\tau$ in the sense of Definition~\ref{def:anti-sigma-derivation} are exactly the Lie derivations of $A$.
\item $\Der(A) \subseteq \LDer(A)$.
\item Let $d \in \Der(A)$, and let $t \colon A \to Z(A)$ be linear with $t([x,y]) = 0$ for all $x, y$. Then $d + t \in \LDer(A)$.
\item In (iii), if $t(1) \neq 0$, then $d + t \notin \Der(A)$.
\end{enumerate}
\end{proposition}

\begin{proof}
We check the four parts in order.
(i) For the swap, $x \cdot_\tau^- y = [x,y]$, so the two identities are the same.
(ii) For $D \in \Der(A)$, $D([x,y]) = D(x)y + xD(y) - D(y)x - yD(x) = [D(x), y] + [x, D(y)]$.
(iii) By (ii), $d([x,y]) = [d(x),y] + [x,d(y)]$, and $t([x,y]) = 0$.
Since $t(x)$ and $t(y)$ are central, $[t(x), y] = 0 = [x, t(y)]$.
Adding gives the Lie derivation identity for $d + t$.
(iv) Every derivation vanishes at $1$, as recalled in Section~\ref{sec:preliminaries}, while $(d+t)(1) = t(1) \neq 0$.
\end{proof}

We now twist the swap as in Section~\ref{sec:twisted}.
The twisted antisymmetrized product is $x \cdot_{\sigma_\alpha}^- y = xy - \alpha(y)\alpha(x)$.

\begin{proposition}\label{prop:twisted-lie}
Let $A$ be a unital associative $K$-algebra, let $\alpha$ be an involutive automorphism or anti-automorphism of $A$, and let $D \colon A \to A$ be a linear map with $D\alpha = \alpha D$.
\begin{enumerate}[label=(\roman*)]
\item If $\alpha$ is an automorphism, then $D$ is a braided-antisymmetric derivation for $\sigma_\alpha$ if and only if $E_D(x,y) - \alpha(E_D(y,x)) = 0$ for all $x, y \in A$.
\item If $\alpha$ is an anti-automorphism, then $D$ is a braided-antisymmetric derivation for $\sigma_\alpha$ if and only if $E_D(x,y) - \alpha(E_D(x,y)) = 0$ for all $x, y \in A$.
\end{enumerate}
\end{proposition}

\begin{proof}
Write $E = E_D$.
For (i), the product is $x \cdot_{\sigma_\alpha}^- y = xy - \alpha(yx)$.
Since $D\alpha = \alpha D$, the left side of the identity is $D(xy) - \alpha(D(yx))$.
The right side is $D(x)y - \alpha(yD(x)) + xD(y) - \alpha(D(y)x)$.
The terms without $\alpha$ give $E(x,y)$ in the difference.
Since $\alpha$ is linear, the terms with $\alpha$ give $-\alpha\bigl(D(yx) - yD(x) - D(y)x\bigr) = -\alpha(E(y,x))$.
So the left side minus the right side is $E(x,y) - \alpha(E(y,x))$.
For (ii), the product is $x \cdot_{\sigma_\alpha}^- y = xy - \alpha(xy)$.
The same steps give the difference $E(x,y) - \alpha(E(x,y))$.
\end{proof}

Take $\alpha = \mathrm{id}$ in (i).
Then the condition says that $E_D$ is symmetric.
This is the condition for a Lie derivation from Section~\ref{sec:twisted}.

Now let $2$ be invertible. % claim-ok: a hypothesis for the next sentences
In case (ii), the condition says that $E_D$ takes values in $A_+$. % claim-ok: the equation E - alpha(E) = 0 says alpha(E) = E, the definition of A_+
In case (i), write $E_D = E_{\mathrm{s}} + E_{\mathrm{a}}$ as in Section~\ref{sec:twisted}.
The condition says that $E_{\mathrm{s}}$ takes values in $A_+$ and $E_{\mathrm{a}}$ takes values in $A_-$.
To see this, swap $x$ and $y$ in (i), then add the two equations and subtract them.
The parities are the opposite of those for the symmetrized product in Proposition~\ref{prop:defect}.

So inner derivations lie in both spaces.
Each $\operatorname{ad}_g$ is a Lie derivation by Proposition~\ref{prop:lie-classification}(ii), and a Jordan derivation by Remark~\ref{rem:sigma-derivation-classification}.
Full matrix algebras show that part (iii) gives maps outside $\Der(A)$.
We use one result on Jordan derivations of matrix rings.
It is stated for a unital associative ring and a $2$-torsion free bimodule \citep[\S2]{Alizadeh2009}.

\begin{theorem}[{\citep[Theorem 3.1]{Alizadeh2009}}]\label{thm:alizadeh}
Let $A$ be a unital associative ring, let $M$ be a $2$-torsion free $A$-bimodule, and let $n \geq 2$. Then every Jordan derivation $D \colon M_n(A) \to M_n(M)$ is a derivation.
\end{theorem}

We also use a description of the first Hochschild cohomology.
It compares the twisted outer derivations of Remark~\ref{rem:inner-twisted} with the ordinary outer derivations, and it gives $HH^1$ of the matrix algebra in Example~\ref{ex:matrix}.
Let $R$ be a $k$-algebra and $M$ an $R$-bimodule.
\citet[\S9.2]{Weibel1994} writes $\mathrm{Der}_k(R,M)$ for the $k$-derivations.
Weibel writes $\mathrm{PDer}_k(R,M)$ for the principal ones $f_m(r) = rm - mr$.

\begin{lemma}[{\citep[Lemma 9.2.1]{Weibel1994}}]\label{lem:weibel}
Let $R$ be a $k$-algebra and $M$ an $R$-bimodule. Then $HH^1(R, M) \cong \mathrm{Der}_k(R,M)/\mathrm{PDer}_k(R,M)$.
\end{lemma}

A nonzero ring $R$ is \emph{simple} if $(0)$ and $R$ are the only ideals in $R$, and a commutative ring is simple if and only if it is a field \citep[\S1]{Lam2001}.

The next theorem shows that $M_n(K)$ is simple.
Example~\ref{ex:matrix} needs this twice: to apply the theorem of Martindale, and to see that $M_n(K)$ is separable.

\begin{theorem}[{\citep[Theorem 3.1]{Lam2001}}]\label{thm:lam-matrix}
Let $R$ be a ring and $M_n(R)$ be the ring of $n \times n$ matrices over $R$. Then any ideal $I$ of $M_n(R)$ has the form $M_n(\mathfrak{m})$ for a uniquely determined ideal $\mathfrak{m}$ of $R$. In particular, if $R$ is a simple ring, so is $M_n(R)$.
\end{theorem}

\begin{theorem}[{\citep[Theorem 2]{Martindale1964}}]\label{thm:martindale}
Let $D$ be a Lie derivation of a primitive ring $R$ into itself, where $R$ contains a nontrivial idempotent and the characteristic of $R$ is not $2$. Then $D$ is of the form $d + t$, where $d$ is an ordinary derivation of $R$ into a primitive ring $\hat{R}$ containing $R$ and $t$ is an additive mapping of $R$ into the center of $\hat{R}$ that maps commutators into zero.
\end{theorem}

Martindale remarks after Theorem~\ref{thm:martindale} that, if $R$ is simple, then $d$ maps $R$ into itself and $t$ maps $R$ into the center of $R$ \citep[p.~187]{Martindale1964}.

\begin{theorem}[{\citep[Theorem 4.1]{Hochschild1945}}]\label{thm:hochschild}
A necessary and sufficient condition for an algebra $\mathfrak{A}$ to be separable is that $H^1(\mathfrak{A}, \mathfrak{M}) = \{0\}$ for every two-sided $\mathfrak{A}$-module $\mathfrak{M}$.
\end{theorem}

We now apply Theorems~\ref{thm:alizadeh}, \ref{thm:lam-matrix}, \ref{thm:martindale} and~\ref{thm:hochschild} to $M_n(K)$.

\begin{example}[Full matrix algebras]\label{ex:matrix}
Let $A = M_n(K)$ with $n \geq 2$, and let $e_{ij}$ be the matrix unit with entry $1$ in position $(i,j)$ and $0$ elsewhere.
If $\operatorname{char} K \neq 2$, every Jordan derivation of $M_n(K)$ is a derivation by Theorem~\ref{thm:alizadeh} applied to the ring $K$ and the $K$-bimodule $K$, since $K$ is then $2$-torsion free.
So $\BSDer_\tau(A) = \JDer(A) = \Der(A)$ in that case.

Let $\theta(x) = \operatorname{tr}(x) I$.
Its values are central, and $\theta([x,y]) = (\operatorname{tr}(xy) - \operatorname{tr}(yx)) I = 0$.
By Proposition~\ref{prop:lie-classification}(iii), $\operatorname{ad}_g + c\,\theta$ is a Lie derivation for every $g \in A$ and $c \in K$.
Since $\theta(e_{11}) = I$ while $\theta(e_{11})e_{11} + e_{11}\theta(e_{11}) = 2e_{11}$ and $\theta(e_{11}^2) = \theta(e_{11})$, the map $\theta$ is not a derivation in any characteristic.

Now let $n \geq 2$ and $\operatorname{char} K \neq 2$.
The ring $M_n(K)$ is simple by Theorem~\ref{thm:lam-matrix}, since a field is simple.
It is primitive, because $K^n$ is a faithful simple left module over it, and $e_{11}$ is a nontrivial idempotent.
By Theorem~\ref{thm:martindale} and the remark after it, every Lie derivation $D$ of $M_n(K)$ has the form $D = d + t$, where $d$ is an additive derivation of $M_n(K)$ into itself and $t$ is an additive map into $K \cdot I$ that vanishes on commutators.

The maps $d$ and $t$ are $K$-linear.
Indeed, for $\lambda \in K$ we have $d(\lambda x) = d(\lambda I)x + \lambda d(x)$, and $d(\lambda I)$ is central because $[d(\lambda I), x] = d([\lambda I, x]) - [\lambda I, d(x)] = 0$.
So $d(\lambda I) = \nu(\lambda) I$ for some $\nu(\lambda) \in K$.
Since $D$ is $K$-linear, $\nu(\lambda)x = \lambda t(x) - t(\lambda x)$ lies in $K \cdot I$ for every $x$.
Taking $x = e_{11}$ gives $\nu(\lambda) = 0$.

Next, $M_n(K)$ is separable, since $M_n(K) \otimes_K F = M_n(F)$ is simple for every extension field $F$ by Theorem~\ref{thm:lam-matrix}.
By Theorem~\ref{thm:hochschild}, $H^1(M_n(K), M_n(K)) = 0$, so every $K$-linear derivation is inner, and $d$ is inner.
Hence $\Der(A) = \Inn(A)$, and every Lie derivation of $A$ has the form $\operatorname{ad}_g + t$ with $t \colon A \to K \cdot I$ linear and $t([A,A]) = 0$.
Such a map $t$ is the same as a $K$-linear map $A/[A,A] \to K$, and $A/[A,A] = HH_0(A,A)$ \citep[\S9.1]{Weibel1994}.
A map $t$ of this kind that is also a derivation is zero.
Indeed $t(e_{ii}) = t(e_{ii}^2) = 2t(e_{ii})e_{ii}$ forces $t(e_{ii}) = 0$ when $n \geq 2$, and $t(e_{ij}) = 0$ for $i \neq j$ because $e_{ij} = [e_{ii}, e_{ij}]$.
So $t$ is determined by $\operatorname{ad}_g + t$: if $\operatorname{ad}_g + t = \operatorname{ad}_h + t'$, then $t - t' = \operatorname{ad}_{h-g}$ is such a derivation, hence zero.
Let $\omega \colon \LDer(A) \to \mathrm{Hom}_K(HH_0(A,A), K)$ send $\operatorname{ad}_g + t$ to $t$.
The map $\omega$ is surjective, since each $t$ is a Lie derivation by Proposition~\ref{prop:lie-classification}(iii), and its kernel is $\Inn(A) = \Der(A)$.
So the sequence
\[
0 \to \Der(A) \to \LDer(A) \xrightarrow{\ \omega\ } \mathrm{Hom}_K\bigl(HH_0(A,A), K\bigr) \to 0
\]
is exact.
The map $\theta$ above is one element of the right-hand space.
\end{example}

The swap gives $\JDer(A)$ through $\mu + \mu \circ \tau$ and $\LDer(A)$ through $\mu - \mu \circ \tau$.
So $\mu$ is half the sum of these two products when $2$ is invertible.
So a map that is both a Jordan and a Lie derivation is a derivation.
On $M_n(K)$ this gives $\JDer(A) \cap \LDer(A) = \Der(A)$ even though $\LDer(A)$ is larger.

For triangular algebras the Lie side has a central term.

A Lie derivation is \emph{proper} when it is the sum of a derivation and a linear map into the centre \citep[\S1]{Cheung2003}.
The extra map vanishes on commutators, so a proper Lie derivation has the standard form of Section~\ref{sec:intro}.

The next theorem uses the algebra $\mathfrak{A} = \mathrm{Tri}(A,M,B)$, and the projections $\pi_A$ and $\pi_B$ onto the diagonal blocks.

The algebra $W(X)$ is the smallest subalgebra of $X$ that satisfies seven closure conditions of \citet[Proposition 10]{Cheung2003}.
We list them for a subalgebra $V$ of $X$.
Here $F$ is a polynomial in one variable over the base ring, and $F'$ is its derivative.
\begin{enumerate}[label=(\roman*)]
\item $V$ contains all commutators $[x, y]$ with $x, y \in X$.
\item If $b \in X$ and $F'(b) = 0$, then $F(b) \in V$.
\item If $b \in X$, $F(b) \in V$ and $F'(b)$ is invertible, then $b \in V$.
\item $V$ contains all idempotents of $X$.
\item $V$ contains all elements $x^p$, where $x \in X$ and $p \geq 0$ is the characteristic of $X$.
\item If $a \in X$ has an inverse $a^{-1}$ in $X$ and $a^{-1} \in V$, then $a \in V$.
\item If $b \in X$ is invertible and $b^k \in V$ for some positive integer $k$, then $kb \in V$.
\end{enumerate}

\begin{theorem}[{\citep[Theorem 11]{Cheung2003}}]\label{thm:cheung}
Let $\mathfrak{A} = \mathrm{Tri}(A, M, B)$. If the following two conditions hold, then every Lie derivation of $\mathfrak{A}$ is proper.
\begin{enumerate}[label=(\Roman*)]
\item $M$ is a faithful left $A$-module and $Z(B) = \pi_B(Z(\mathfrak{A}))$; or $M$ is a faithful left $A$-module and $A = W(A)$; or every Lie derivation of $A$ is proper and $A = W(A)$.
\item $M$ is a faithful right $B$-module and $Z(A) = \pi_A(Z(\mathfrak{A}))$; or $M$ is a faithful right $B$-module and $B = W(B)$; or every Lie derivation of $B$ is proper and $B = W(B)$.
\end{enumerate}
\end{theorem}

\citet[\S3]{LiWei2015} extend this analysis to generalized matrix algebras.
A generalized matrix algebra \citep[\S2.1]{LiWei2015} is built from two algebras $A$ and $B$ and two bimodules ${}_A M_B$ and ${}_B N_A$.
It also has two bimodule maps
\[
M \otimes_B N \to A, \qquad N \otimes_A M \to B.
\]
These data form a Morita context \citep[\S2.1]{LiWei2015}.
The bimodule $N$ may be zero, and then the algebra is triangular \citep[\S2.1 and Corollary 3.6]{LiWei2015}.

\begin{corollary}\label{cor:triangular-lie}
Let $T = \mathrm{Tri}(A,M,B)$ satisfy conditions (I) and (II) of Theorem~\ref{thm:cheung}. Then every braided-antisymmetric derivation of $T$ for $\tau$ has the standard form $d + t$, and it is a derivation only if $t(1) = 0$.
\end{corollary}

\begin{proof}
By Proposition~\ref{prop:lie-classification}(i), the maps in question are the Lie derivations of $T$, which are proper by Theorem~\ref{thm:cheung}. The last claim is Proposition~\ref{prop:lie-classification}(iv).
\end{proof}

Corollary~\ref{cor:triangular-jordan} says that on such a triangular algebra the symmetrized product gives only derivations.
Example~\ref{ex:t2} shows that the antisymmetrized product can add a central term.
Take $T_2(K) = \mathrm{Tri}(K,K,K)$ over a field with $\operatorname{char} K \neq 2$.
Here $M = K$ is faithful on both sides.
Also $Z(T_2(K)) = KI$ projects onto $K$.
So Corollaries~\ref{cor:triangular-jordan} and~\ref{cor:triangular-lie} apply.

\begin{example}[Upper triangular $2 \times 2$ matrices]\label{ex:t2}
Let $T = T_2(\mathbb{F}_3)$.
Every additive map $T \to T$ is $\mathbb{F}_3$-linear.
So it is one of the $3^9$ maps given by a $3 \times 3$ matrix over $\mathbb{F}_3$.

An exhaustive computer check over all of them confirms Corollaries~\ref{cor:triangular-jordan} and~\ref{cor:triangular-lie}.
Every Jordan derivation is a derivation.
Every Lie derivation equals $\operatorname{ad}_x + t$ for some $x \in T$ and some central-valued $t$ that vanishes on commutators.

The term $t$ is needed.
Define $\theta$ by
\[
\theta\begin{pmatrix} a & m \\ 0 & b \end{pmatrix} = aI.
\]
It has central values and vanishes on commutators, whose diagonal is zero.
So it is a Lie derivation.
It is not a derivation, since
\[
\theta(e_1^2) = I \qquad \text{while} \qquad \theta(e_1)e_1 + e_1 \theta(e_1) = 2e_1.
\]
So on $T_2(\mathbb{F}_3)$, as on $M_n(K)$, we have $\JDer = \Der$, while $\LDer$ is strictly larger.
\end{example}

\begin{thebibliography}{9}

\bibitem[Alizadeh(2009)]{Alizadeh2009}
Alizadeh, R.: Jordan derivations of full matrix algebras. Linear Algebra Appl. \textbf{430}, 574--578 (2009).

\bibitem[Amitsur and Levitzki(1950)]{Amitsur1950}
Amitsur, S.A., Levitzki, J.: Minimal identities for algebras. Proc. Amer. Math. Soc. \textbf{1}(4), 449--463 (1950).

\bibitem[Benkovi\v{c}(2015)]{Benkovic2015}
Benkovi\v{c}, D.: A note on $f$-derivations of triangular algebras. Aequationes Math. \textbf{89}, 1207--1211 (2015).


\bibitem[Bourbaki(1989)]{Bourbaki1989}
Bourbaki, N.: Algebra I. Chapters 1--3. Elements of Mathematics. Springer, Berlin (1989).

\bibitem[Bre\v{s}ar(1988)]{Bresar1988}
Bre\v{s}ar, M.: Jordan derivations on semiprime rings. Proc. Amer. Math. Soc. \textbf{104}(4), 1003--1006 (1988).

\bibitem[Bre\v{s}ar(1993)]{Bresar1993}
Bre\v{s}ar, M.: Commuting traces of biadditive mappings, commutativity-preserving mappings and Lie mappings. Trans. Amer. Math. Soc. \textbf{335}(2), 525--546 (1993).

\bibitem[Chambert-Loir(2021)]{ChambertLoir2021}
Chambert-Loir, A.: (Mostly) Commutative Algebra. Universitext. Springer, Cham (2021).

\bibitem[Cheung(2003)]{Cheung2003}
Cheung, W.-S.: Lie derivations of triangular algebras. Linear Multilinear Algebra \textbf{51}(3), 299--310 (2003).

\bibitem[Cusack(1975)]{Cusack1975}
Cusack, J.M.: Jordan derivations on rings. Proc. Amer. Math. Soc. \textbf{53}(2), 321--324 (1975).

\bibitem[Herstein(1957)]{Herstein1957}
Herstein, I.N.: Jordan derivations of prime rings. Proc. Amer. Math. Soc. \textbf{8}(6), 1104--1110 (1957).

\bibitem[Herstein(1976)]{Herstein1976}
Herstein, I.N.: Rings with Involution. Chicago Lectures in Mathematics. University of Chicago Press, Chicago (1976).

\bibitem[Hochschild(1945)]{Hochschild1945}
Hochschild, G.: On the cohomology groups of an associative algebra. Ann. of Math. \textbf{46}(1), 58--67 (1945).

\bibitem[Hosseini and Karizaki(2019)]{HosseiniKarizaki2019}
Hosseini, A., Karizaki, M.M.: Functional equations characterizing $\sigma$-derivations. Acta Sci. Math. (Szeged) \textbf{85}, 431--440 (2019).

\bibitem[Jacobson(1937)]{Jacobson1937}
Jacobson, N.: Abstract derivation and Lie algebras. Trans. Amer. Math. Soc. \textbf{42}(2), 206--224 (1937).

\bibitem[Huru(2006)]{Huru2006}
Huru, H.L.: Quantizations of braided derivations. I: Monoidal categories. Lobachevskii J. Math. \textbf{24}, 13--42 (2006).

\bibitem[Li and Wei(2015)]{LiWei2015}
Li, Y.-B., Wei, F.: Lie derivations of dual extensions of algebras. Colloq. Math. \textbf{141}(1), 65--82 (2015). arXiv:1303.0831.

\bibitem[Lam(2001)]{Lam2001}
Lam, T.Y.: A First Course in Noncommutative Rings, 2nd edn. Graduate Texts in Mathematics 131. Springer, New York (2001).

\bibitem[Martindale(1964)]{Martindale1964}
Martindale, W.S., III: Lie derivations of primitive rings. Michigan Math. J. \textbf{11}(2), 183--187 (1964).

\bibitem[Rotman(2009)]{Rotman2009}
Rotman, J.J.: An Introduction to Homological Algebra, 2nd edn. Universitext, Springer, New York (2009).

\bibitem[R\'o\.za\'nski(1996)]{Rozanski1996}
R\'o\.za\'nski, J.: Bicharacters, braids and Jacobi identity. arXiv:q-alg/9611029 (1996).

\bibitem[Weber(2020)]{Weber2020}
Weber, T.: Braided Cartan calculi and submanifold algebras. J. Geom. Phys. \textbf{156}, 103612 (2020). \url{https://doi.org/10.1016/j.geomphys.2020.103612}

\bibitem[Zhang and Yu(2006)]{ZhangYu2006}
Zhang, J.-H., Yu, W.-Y.: Jordan derivations of triangular algebras. Linear Algebra Appl. \textbf{419}(1), 251--255 (2006).

\bibitem[Wang(2019)]{Wang2019}
Wang, Y.: Lie (Jordan) derivations of arbitrary triangular algebras. Aequationes Math. \textbf{93}, 1221--1229 (2019).

\bibitem[Weibel(1994)]{Weibel1994}
Weibel, C.A.: An Introduction to Homological Algebra. Cambridge Stud. Adv. Math. 38, Cambridge University Press, Cambridge (1994).

\end{thebibliography}

\end{document}
